\section{对数螺线的等角性}
\markboth{92. 对数螺线的等角性}{92. 对数螺线的等角性}

对数螺线的极坐标方程是
\[
  \rho = e^{m\theta},
\]
其中 $m$ 为常数，并且 $m \neq 0$。对 $\theta$ 求导，得
\[
  \rho' = m e^{m\theta},
\]
因而
\[
  \frac{\rho}{\rho'} = \frac{e^{m\theta}}{m e^{m\theta}} = \frac{1}{m} = \text{常数}.
\]

设对数螺线在点 $M(\rho, \theta)$ 的切线与极半径 $OM$ 的夹角为 $\psi$，根据第 51 节的公式，得到
\[
  \tg\psi = \frac{\rho}{\rho'} = \frac{1}{m},
\]
因而对于对数螺线，角 $\psi$ 保持常数。这个性质叫做对数螺线的等角性，并且对数螺线由于这个特性而得到“等角螺线”的别名。

在金属切削加工时，刀具与金属工件之间存在着很大的切削力，所以对于刀刃的前后两面与工件表面的夹角都有一定的要求。特别是一种叫做成型铣刀的旋转刀具，磨刀时只磨刀齿的前面，不磨刀齿的背面，因而对齿背曲线要求更高，必须使它在每一点的切线都与过该点的半径成定角。对数螺线正好能满足这个要求。所以，成型铣刀的刀齿背面总是做成对数螺线的形状。
